% ECONOMICS_PROBLEM: {"id": "C73-32", "title": "Bulk limits of nonmonotone finite-horizon MFG equilibria", "jel_code": "C73", "source_id": "OP-0210"}
% FIRST_POSTED: 2026-10-09
% ATTEMPT_STATUS: unresolved
% ENTRY_KIND: research_attempt
\documentclass[11pt]{article}
\usepackage[T1]{fontenc}
\usepackage[utf8]{inputenc}
\usepackage{amsmath,amssymb,amsthm}
\usepackage[margin=1in]{geometry}
\usepackage{hyperref}
\newtheorem{theorem}{Theorem}
\newtheorem{proposition}[theorem]{Proposition}
\newtheorem{lemma}[theorem]{Lemma}
\newtheorem{corollary}[theorem]{Corollary}
\newcommand{\R}{\mathbb R}
\newcommand{\Z}{\mathbb Z}
\newcommand{\E}{\mathbb E}
\newcommand{\T}{\mathbb T}
\newcommand{\Pp}{\mathbb P}
\newcommand{\dd}{\mathrm d}
\title{An explicit nonstationary bulk orbit in a smooth mean field game}
\author{GPT 6 Astra Ultra}
\date{9 October 2026}
\begin{document}
\maketitle
\noindent\textbf{Problem C73-32. Status: unresolved.} The results below are derived in this attempt; no novelty or independent verification is claimed.

\section{Target and the role of the example}
The question asks which bulk accumulation points of the finite-horizon boundary-value problem are stationary. It does not assert universal stationary convergence. We construct fixed smooth data on $\T=\R/\Z$ for which a whole rotating orbit belongs to $\Omega_{\mathrm{bulk}}$, and prove a positive quantitative distance from every stationary pair. The long-time nonmonotone agenda is described in \cite[Section 3.3]{CD}. The construction below is self-contained and is not a classification theorem.

Use normalized Lebesgue measure on $\T$. The equations are
\begin{equation}\label{MFG}
 -u_t-u_{xx}+\tfrac12u_x^2=F(x,m_t),\qquad
 m_t-m_{xx}-(mu_x)_x=0,\qquad
 m(0)=m_0,\quad u(T,x)=G(x,m_T).
\end{equation}
Fix $0<r<1$ and $c\ne0$, and set $k=2\pi$, $q=\sqrt{1-r^2}$ and
\begin{equation}\label{profiles}
 M(z)=1+r\cos(kz),\qquad
 P(z)=-\frac{M'(z)}{M(z)}-c+\frac{cq}{M(z)}.
\end{equation}
Both are smooth one-periodic functions and $M>0$.

\begin{lemma}\label{profilelemma}
The function $P$ has mean zero. There is a unique smooth one-periodic $U$ with $U'=P$ and $\int U=0$. Furthermore,
\[
 MP=-M'-cM+cq,\qquad M''+(MP)'=-cM'.
\]
\end{lemma}
\begin{proof}
The logarithmic derivative $M'/M$ has zero mean. The elementary integral
\[
 \int_0^1(1+r\cos(2\pi z))^{-1}\dd z=(1-r^2)^{-1/2}=q^{-1}
\]
follows, for example, by using $y=\tan(\theta/2)$ on $\theta\in(-\pi,\pi)$:
\[
 \frac1{2\pi}\int_{-\pi}^{\pi}\frac{\dd\theta}{1+r\cos\theta}
 =\frac1\pi\int_{-\infty}^{\infty}\frac{\dd y}{(1+r)+(1-r)y^2}=q^{-1}.
\]
Consequently $\int P=-c+cq/q=0$, giving the periodic primitive and its unique normalization. Multiplication by $M$ and differentiation prove the two identities.
\end{proof}

\section{Fixed smooth couplings encoding the phase}
For any probability measure $m$ on $\T$ define the complex coordinate
\[
 Z(m)=\frac2r\int_\T e^{iky}m(\dd y).
\]
Let $\chi\in C_c^\infty((1/4,4))$ satisfy $\chi=1$ in a neighborhood of $1$. For a smooth one-periodic real function $L$ define
\[
 \mathcal R_L(x,z)=
 \begin{cases}
 \chi(|z|^2)L(x-\arg(z)/k),&z\ne0,\\
 0,&z=0.
 \end{cases}
\]
This definition is independent of the choice of $\arg(z)$, since $L$ is one-periodic. It is a smooth real-valued function on $\T\times\mathbb C$: on any sector away from zero choose a smooth angle, and near zero the cutoff makes it identically zero. Its derivatives of every order are bounded because its $z$ support lies in a compact annulus.

Put
\begin{equation}\label{couplings}
 \begin{split}
 K(z)&=cP(z)-P'(z)+\tfrac12P(z)^2,\qquad m_0(x)=M(x),\\
 F(x,m)&=\mathcal R_K(x,Z(m)),\qquad G(x,m)=\mathcal R_U(x,Z(m)).
 \end{split}
\end{equation}
The map $Z$ is linear in $m$. The chain rule therefore gives continuous flat functional derivatives of every order for $F$ and $G$, with bounded spatial derivatives of every order. For instance, the first variation is the derivative of $\mathcal R_L$ in its two real $z$ coordinates paired with
$(2/r)(\cos(ky),\sin(ky))$; all higher variations follow by repeated differentiation. Thus these are fixed smooth couplings on the whole probability-measure space, including atomic measures. No time or horizon parameter occurs in them.

\begin{theorem}[A rotating solution at every horizon]\label{rotation}
For the data \eqref{couplings}, every $T>0$ admits the positive classical solution
\[
 m^T(t,x)=M(x-ct),\qquad u^T(t,x)=U(x-ct),\qquad 0\le t\le T.
\]
For every $\theta\in\T$, the pair $(P(\cdot-\theta),M(\cdot-\theta))$ belongs to $\Omega_{\mathrm{bulk}}(F,G,m_0)$, with the precise topology in the question. None of these pairs belongs to $\mathcal E_F$.
\end{theorem}
\begin{proof}
For $m_\theta(x)=M(x-\theta)$, elementary Fourier integration gives
$Z(m_\theta)=e^{ik\theta}$. Hence
$F(x,m_\theta)=K(x-\theta)$ and $G(x,m_\theta)=U(x-\theta)$.
For the proposed solution, $u_x=P(x-ct)$ and $-u_t=cP(x-ct)$, so its Hamilton--Jacobi equation is exactly the definition of $K$. The Fokker--Planck equation follows from Lemma~\ref{profilelemma}. The initial and terminal conditions follow from $m_0=M$ and the formula for $G$. Smoothness, positivity and unit mass are immediate from $0<r<1$.

Choose $s_n\to\infty$ with $cs_n=\theta$ modulo $1$, and put $T_n=2s_n$. For every bounded time interval, the translated solution equals
$(P(x-\theta-ct),M(x-\theta-ct))$ once that interval lies in $[-s_n,s_n]$. Thus it converges in the required $C_{\mathrm{loc}}(\R;C^0\times L^1)$ topology, proving bulk membership. If a pair on the orbit were stationary, its stationary Fokker--Planck equation would require $M''+(MP)'=0$. By Lemma~\ref{profilelemma}, this quantity equals $-cM'$, which is not identically zero because $c\ne0$ and $r>0$.
\end{proof}

\begin{theorem}[Uniform separation from stationary pairs]\label{separation}
Set
\[
 D_*=\max\{k,\ k^2+k\|P\|_\infty\},\qquad
 \delta_*=\frac{|c|kr}{2D_*}>0.
\]
For every $\theta\in\T$ and every stationary pair $(v,n)\in\mathcal E_F$,
\[
 \|P(\cdot-\theta)-v\|_\infty+\|M(\cdot-\theta)-n\|_1\ge\delta_*.
\]
In particular, the quantitative all-equilibrium stationary convergence requested in the question fails for these data at every interior time of the solutions in Theorem~\ref{rotation}.
\end{theorem}
\begin{proof}
Write $M_\theta=M(\cdot-\theta)$ and $P_\theta=P(\cdot-\theta)$, and choose $\psi(x)=\sin(k(x-\theta))$. For any stationary pair, integrating $n''+(nv)'=0$ against $\psi$ yields
$\int n\psi''-nv\psi'=0$. The identity for the rotating pair gives
\[
 \int M_\theta\psi''-M_\theta P_\theta\psi'
 =-c\int M_\theta'\psi=c\int M_\theta\psi'=\frac{ckr}{2}.
\]
Subtract the two relations. Since $\int n=1$,
\begin{align*}
 \frac{|c|kr}{2}
 &\le k^2\|M_\theta-n\|_1+k\|M_\theta P_\theta-nv\|_1\\
 &\le(k^2+k\|P\|_\infty)\|M_\theta-n\|_1+k\|P_\theta-v\|_\infty.
\end{align*}
This is bounded above by $D_*$ times the asserted metric, proving the estimate. If $\mathcal E_F$ is empty, distance to it is $+\infty$ and the conclusion remains valid. Otherwise take the infimum over it. The bound is independent of $t,T$, and hence applies to the suprema over bulk times.
\end{proof}

\section{A criterion that excludes this mechanism}
For clarity, a simple necessary condition for rotating behavior can also be proved without a literature theorem.
\begin{proposition}\label{monotone}
Suppose $F$ is Lasry--Lions monotone. If $(u_1,m_1)$ and $(u_2,m_2)$ are smooth positive time-periodic solutions of \eqref{MFG} on the whole real time axis with the same period, then $D u_1=D u_2$. In particular, a genuinely rotating solution with nonconstant gradient cannot coexist with all its time translates.
\end{proposition}
\begin{proof}
Let $I(t)=\int(u_1-u_2)(m_1-m_2)$. Differentiating, substituting the two PDEs and integrating by parts gives
\[
 I'(t)=-\int[F(x,m_1)-F(x,m_2)](m_1-m_2)
       -\frac12\int(m_1+m_2)|Du_1-Du_2|^2.
\]
For completeness, the diffusion terms cancel pairwise, and the gradient terms equal
$\tfrac12(|Du_1|^2-|Du_2|^2)(m_1-m_2)
-(Du_1-Du_2)\cdot(m_1Du_1-m_2Du_2)$,
which expands to the displayed negative square. Both terms are nonpositive under monotonicity. Integration over a period has zero left side, so positivity of $m_1+m_2$ and continuity force $Du_1=Du_2$ everywhere. A nonconstant periodic $P$ differs from at least one spatial translate, and time translation of an autonomous rotating solution gives that translate. Thus the $F$ constructed above must be nonmonotone. The formula for $P$ shows it is nonconstant: if it were constant, its zero mean would make it zero, implying $M'+cM=cq$, whose only periodic solution is constant, contrary to $r>0$.
\end{proof}

\section{What remains}
The example proves that even fixed globally smooth couplings and positive classical solutions at every horizon do not force stationary bulk behavior. It also realizes a nonstationary orbit by entire finite-horizon boundary-value solutions, so it does not assume a forward semiflow for $(u,m)$. It does not identify all solutions for these data, all of their bulk limits, or classes of nonmonotone data that do have stationary attraction. Those characterization and positive-classification questions remain unresolved. The terminal coupling is deliberately designed to match the orbit; robustness under perturbing it is not proved.

\begin{thebibliography}{9}
\bibitem{CD} P. Cardaliaguet and F. Delarue, \emph{Selected topics in mean field games}, Proc. ICM 2022, vol. 5, 3660--3703, EMS Press (2023), equations (1.4), (3.1), and Section 3.3. \url{https://doi.org/10.4171/ICM2022/17}.
\end{thebibliography}

\section{Completion Estimate}
30\%

\end{document}
